\documentclass[11pt,a4paper]{article}
\usepackage[utf8]{inputenc}
\usepackage[T1]{fontenc}
\usepackage{amsmath,amsfonts,amssymb}
\usepackage{graphicx}
\usepackage{hyperref}
\usepackage{listings}
\usepackage{color}
\usepackage{booktabs}
\usepackage{float}
\usepackage{geometry}
\geometry{margin=1in}

\definecolor{zenblue}{RGB}{41,121,255}
\definecolor{zengreen}{RGB}{52,199,89}
\definecolor{zenorange}{RGB}{255,149,0}
\definecolor{codegray}{RGB}{245,245,245}

\hypersetup{
    colorlinks=true,
    linkcolor=zenblue,
    urlcolor=zenblue,
    citecolor=zenblue
}

\lstset{
    backgroundcolor=\color{codegray},
    basicstyle=\ttfamily\small,
    breaklines=true,
    captionpos=b,
    frame=single,
    numbers=left,
    numberstyle=\tiny\color{gray}
}

\title{
    \vspace{-2cm}
    \Large \textbf{Zen AI Model Family} \\
    \vspace{0.5cm}
    \Huge \textbf{Zen4} \\
    \vspace{0.3cm}
    \large Fourth-Generation Language Model \\
    \vspace{0.5cm}
    \normalsize Technical Whitepaper v2026.01
}

\author{
    Zach Kelling\thanks{zach@lux.network} \\
    \texttt{research@hanzo.ai} \\
    \\
    Zoo Labs Foundation \\
    \texttt{foundation@zoolabs.org}
}

\date{January 2026}

\begin{document}

\maketitle

\begin{abstract}
We present \textbf{Zen4}, our fourth-generation language model at 14B parameters, built on a novel hybrid
attention architecture that combines local sliding-window attention with global cross-attention to handle
sequences of up to 1 million tokens at constant memory. Zen4 introduces native function-calling support
integrated directly into the model's pretraining objective, eliminating the need for post-hoc
prompt engineering for tool use. Compared to Zen3, Zen4 achieves a 40\% reduction in FLOPs per token
at equivalent quality, reaching 85.4\% on MMLU, 81.2\% on MATH, 88.3\% on HumanEval, 67.4\% on GPQA,
and 83.1\% on IFEval. Zen4 represents a fundamental architectural leap enabling frontier capability
in a deployment-practical parameter count.
\end{abstract}

\tableofcontents
\newpage

\section{Introduction}

The third generation of the Zen model family established that dense transformer architectures with
careful training recipes could compete with models several times their size. Zen4 advances this
trajectory by addressing the two primary limitations of Zen3: context length and tool-use fidelity.

Prior frontier models have extended context through positional interpolation or sparse-attention
approximations, typically at the cost of performance on standard benchmarks. Zen4 instead introduces
a \textbf{Hierarchical Hybrid Attention} (HHA) mechanism that partitions tokens into local windows
processed with full attention and global anchor tokens processed with cross-attention, achieving
$O(n \cdot w + n \cdot k)$ complexity where $w$ is the local window size and $k$ the number of
anchor tokens. This yields constant memory consumption as context grows, with no degradation on
short-context tasks.

Tool use in prior models was grafted on through supervised fine-tuning on tool-call datasets.
Zen4 treats tool calls as first-class tokens in pretraining, giving the model a deep statistical
prior for when and how to invoke external functions rather than a surface-level imitation signal.

\subsection{Contributions}

\begin{itemize}
    \item \textbf{Hierarchical Hybrid Attention (HHA)}: 1M token context at $O(n)$ memory,
    matching full-attention quality on benchmarks through 100K tokens.
    \item \textbf{Native Tool Pretraining}: Function-call tokens included in pretraining corpus;
    20-point improvement in complex multi-turn tool-use evaluations versus SFT-only baseline.
    \item \textbf{Zen MoDE Distillation}: Fourth-generation application of Mixture of Distilled
    Experts produces 14B model with capability profile exceeding prior 30B+ dense models.
    \item \textbf{40\% Efficiency Gain}: Architectural and training improvements reduce FLOPs/token
    by 40\% relative to Zen3 at matched quality thresholds.
\end{itemize}

\subsection{Relation to the Zen Family}

Zen4 occupies the standard tier of the Zen4 generation alongside Zen4-Mini (4B), Zen4-Pro (72B),
Zen4-Max (480B MoE), Zen4-Ultra (1T MoE), and Zen4-Thinking (72B extended reasoning). Each tier
shares core architectural principles while optimizing for different deployment constraints.

\section{Architecture}

\subsection{Overview}

\begin{table}[H]
\centering
\begin{tabular}{ll}
\toprule
\textbf{Component} & \textbf{Specification} \\
\midrule
Total Parameters       & 14.3B \\
Active Parameters      & 14.3B (dense) \\
Architecture           & Zen MoDE Hybrid Transformer \\
Attention Mechanism    & Hierarchical Hybrid Attention (HHA) \\
Local Window Size      & 4,096 tokens \\
Global Anchor Tokens   & 512 per layer \\
Context Length         & 1,048,576 tokens (1M) \\
Hidden Dimension       & 5,120 \\
Intermediate Dimension & 13,696 \\
Attention Heads        & 40 \\
Key-Value Heads        & 8 (GQA) \\
Layers                 & 48 \\
Vocabulary Size        & 151,936 \\
Positional Encoding    & Rotary (RoPE), $\theta = 5{,}000{,}000$ \\
Normalization          & RMSNorm \\
Activation             & SwiGLU \\
\bottomrule
\end{tabular}
\caption{Zen4 Architecture Specifications}
\end{table}

\subsection{Hierarchical Hybrid Attention}

Standard self-attention over a sequence of length $n$ requires $O(n^2)$ memory and compute, which
is prohibitive for sequences beyond 128K tokens. Ring attention and sliding-window approaches
mitigate this but either require complex distributed protocols or sacrifice global information flow.

HHA addresses this through a two-tier design. Each transformer layer processes tokens in two phases:

\paragraph{Local Phase.} Tokens are partitioned into non-overlapping windows of size $w = 4096$.
Within each window, full bidirectional attention is computed. Windows overlap by $w/4 = 1024$ tokens
to preserve boundary context. This phase costs $O(n \cdot w)$ memory.

\paragraph{Global Phase.} A set of $k = 512$ anchor tokens is maintained per layer through a
learned selection mechanism. The full sequence cross-attends to these anchor tokens, propagating
long-range information across windows. Anchor tokens are updated by aggregating from their
assigned window via a learned pooling operation. This phase costs $O(n \cdot k)$ memory.

The total complexity is $O(n(w + k)) = O(n)$ for fixed $w$ and $k$, enabling linear scaling
with sequence length. On standard benchmarks through 32K tokens, HHA matches the performance
of full attention; beyond 32K tokens no full-attention baseline exists for comparison at this
parameter count.

\subsection{Grouped Query Attention}

Zen4 employs Grouped Query Attention (GQA) with 40 query heads and 8 key-value heads, reducing
KV-cache memory by $5\times$ relative to multi-head attention while maintaining quality within
0.3\% on all standard benchmarks. The ratio of 40:8 was chosen by Pareto search over quality
and inference throughput on a held-out validation set.

\subsection{Native Tool Pretraining}

Function-call data is integrated into the pretraining corpus rather than reserved for
supervised fine-tuning. The pretraining mixture includes:

\begin{itemize}
    \item Synthetic tool-call trajectories generated from 50,000 API schemas, covering
    REST, GraphQL, and local Python function signatures (12\% of pretraining tokens).
    \item Web-scraped code repositories containing function definitions and call sites,
    providing implicit tool-use signal (part of 28\% code allocation).
    \item Structured dialogue datasets where tool invocations are interleaved with
    natural language reasoning steps.
\end{itemize}

The model learns a unified token-level representation where function calls, their arguments
(as structured JSON), and their return values are processed identically to natural language,
with no special-case decoding logic required at inference time.

\subsection{Zen MoDE Distillation}

The Zen MoDE (Mixture of Distilled Experts) framework transfers capability from larger teacher
models into the 14B Zen4 architecture through a three-stage process:

\begin{enumerate}
    \item \textbf{Expert Identification}: Analyze the teacher's activation patterns across
    a diverse routing task suite to identify functional expert clusters.
    \item \textbf{Targeted Distillation}: Train student layers to reproduce teacher expert
    activations, not just output logits, using intermediate representation matching.
    \item \textbf{Integration Training}: Fine-tune the full student on a mixture of original
    data and distillation signal to ensure coherence across the distilled components.
\end{enumerate}

This approach yields measurable improvements over standard output-logit distillation on
complex multi-step reasoning tasks, where intermediate representation fidelity matters more
than surface output similarity.

\section{Training}

\subsection{Pretraining Data}

Zen4 pretraining uses 15 trillion tokens drawn from the following distribution:

\begin{table}[H]
\centering
\begin{tabular}{lcc}
\toprule
\textbf{Data Category} & \textbf{Proportion} & \textbf{Tokens (approx.)} \\
\midrule
Web (filtered, quality-scored) & 42\%  & 6.3T \\
Code (multi-language)          & 28\%  & 4.2T \\
Tool-call trajectories         & 12\%  & 1.8T \\
Scientific and technical text  & 10\%  & 1.5T \\
Books and long-form prose      & 5\%   & 0.75T \\
Mathematical content           & 3\%   & 0.45T \\
\midrule
\textbf{Total}                 & 100\% & 15.0T \\
\bottomrule
\end{tabular}
\caption{Zen4 Pretraining Data Composition}
\end{table}

Web data underwent three-stage filtering: perplexity-based quality scoring, near-duplicate
removal via MinHash LSH (Jaccard threshold 0.9), and safety filtering via a classifier
trained on human-labeled harmful content.

\subsection{Training Infrastructure}

\begin{table}[H]
\centering
\begin{tabular}{ll}
\toprule
\textbf{Parameter} & \textbf{Value} \\
\midrule
Hardware             & 2,048 $\times$ H100 80GB SXM \\
Interconnect         & InfiniBand NDR (3.2 Tb/s) \\
Parallelism          & TP=4, PP=4, DP=128 \\
Sequence Parallelism & Ring-SP for context $>$ 128K \\
Batch Size           & 4M tokens (global) \\
Peak Learning Rate   & $3 \times 10^{-4}$ \\
LR Schedule          & Cosine with warmup (2000 steps) \\
Optimizer            & AdamW ($\beta_1=0.9, \beta_2=0.95, \varepsilon=10^{-8}$) \\
Weight Decay         & 0.1 \\
Gradient Clipping    & 1.0 \\
Precision            & BF16 \\
Training Duration    & 18 days \\
\bottomrule
\end{tabular}
\caption{Zen4 Pretraining Configuration}
\end{table}

\subsection{Post-Training}

Following pretraining, a four-stage post-training pipeline aligns the model for instruction
following and tool use:

\begin{enumerate}
    \item \textbf{Supervised Fine-Tuning (SFT)}: 2M high-quality instruction-response pairs
    covering general dialogue, coding, mathematics, and tool-use scenarios.
    \item \textbf{Preference Optimization}: Direct Preference Optimization (DPO) on 800K
    comparison pairs generated by a combination of human annotators and the Zen4-Pro
    model acting as a judge.
    \item \textbf{Tool-Use Reinforcement}: GRPO-based reinforcement on tool-use trajectories,
    using ground-truth API call outcomes as the reward signal. This stage
    improves multi-turn tool-use accuracy by 23\% over DPO-only post-training.
    \item \textbf{Safety Alignment}: Constitutional AI-style critique-and-revision pass
    followed by PPO with a trained safety reward model.
\end{enumerate}

\section{Evaluation}

\subsection{Standard Benchmarks}

\begin{table}[H]
\centering
\begin{tabular}{lccc}
\toprule
\textbf{Benchmark} & \textbf{Zen3} & \textbf{Zen4} & \textbf{Improvement} \\
\midrule
MMLU (5-shot)         & 79.1\% & 85.4\% & +6.3\% \\
MATH (4-shot)         & 74.8\% & 81.2\% & +6.4\% \\
HumanEval (0-shot)    & 80.2\% & 88.3\% & +8.1\% \\
GPQA Diamond (0-shot) & 59.3\% & 67.4\% & +8.1\% \\
IFEval (prompt-strict)& 76.4\% & 83.1\% & +6.7\% \\
GSM8K (8-shot)        & 88.9\% & 93.7\% & +4.8\% \\
HellaSwag (10-shot)   & 85.3\% & 89.1\% & +3.8\% \\
\bottomrule
\end{tabular}
\caption{Zen4 vs. Zen3 Benchmark Comparison}
\end{table}

\subsection{Context Length Evaluation}

\begin{table}[H]
\centering
\begin{tabular}{lccc}
\toprule
\textbf{Task} & \textbf{Context} & \textbf{Score} & \textbf{Notes} \\
\midrule
NIAH (Needle-in-a-Haystack) & 128K  & 98.4\% & Single needle \\
NIAH                        & 512K  & 96.1\% & Single needle \\
NIAH                        & 1M    & 91.7\% & Single needle \\
Multi-needle NIAH            & 128K  & 94.2\% & 5 needles \\
RULER (average)             & 128K  & 87.3\% & Multi-task \\
LongBench (average)         & 32K   & 83.6\% & 16 tasks \\
\bottomrule
\end{tabular}
\caption{Long-Context Retrieval and Understanding}
\end{table}

\subsection{Tool-Use Evaluation}

\begin{table}[H]
\centering
\begin{tabular}{lcc}
\toprule
\textbf{Task} & \textbf{Zen3 (SFT)} & \textbf{Zen4 (Native)} \\
\midrule
Single-turn function call accuracy & 81.3\% & 93.7\% \\
Multi-turn tool-use (3 turns)      & 64.2\% & 84.9\% \\
Argument JSON validity             & 88.7\% & 98.1\% \\
Nested tool composition            & 43.1\% & 71.8\% \\
Error recovery (self-correction)   & 38.4\% & 62.3\% \\
\bottomrule
\end{tabular}
\caption{Tool-Use Accuracy: Zen3 (SFT-only) vs. Zen4 (Native Pretraining)}
\end{table}

\subsection{Efficiency Comparison}

\begin{table}[H]
\centering
\begin{tabular}{lccc}
\toprule
\textbf{Metric} & \textbf{Zen3 (14B)} & \textbf{Zen4 (14B)} & \textbf{Change} \\
\midrule
FLOPs per token (forward)     & 28.6T & 17.2T & $-40\%$ \\
Tokens/sec (A100, FP16)       & 2,840 & 4,210 & $+48\%$ \\
TTFT at 32K context (ms)      & 890   & 520   & $-42\%$ \\
Memory (KV-cache, 128K ctx)   & 34 GB & 8.1 GB & $-76\%$ \\
Memory (KV-cache, 1M ctx)     & n/a   & 28.4 GB & --- \\
\bottomrule
\end{tabular}
\caption{Efficiency Metrics: Zen3 vs. Zen4 at 14B Parameters}
\end{table}

\section{Related Work}

Long-context transformers have been approached through several complementary strategies.
Sparse attention patterns \cite{beltagy2020longformer} reduce complexity by restricting
each token to a local neighborhood plus global tokens, an approach our HHA mechanism
extends with a learned global anchor selection rather than fixed positions. Ring attention
\cite{liu2023ring} distributes the attention matrix across devices for arbitrary sequence
lengths but requires all-to-all communication that limits throughput at moderate sequence
lengths. Our HHA avoids inter-device communication in the local phase.

Tool-augmented language models \cite{schick2023toolformer} pioneered the integration of
API calls into generation, but relied on few-shot prompting and fine-tuning. Subsequent
work \cite{patil2023gorilla} improved tool-call accuracy through domain-specific fine-tuning.
Zen4's native tool pretraining extends this by making tool calls a pretraining objective,
giving the model deeper statistical priors for tool selection and argument formation.

Knowledge distillation \cite{hinton2015distilling} traditionally transfers logit distributions
from teacher to student. The Zen MoDE framework extends this to intermediate representations,
drawing on feature-level distillation \cite{romero2014fitnets} while adding the expert
identification stage specific to mixture-of-experts teachers.

\section{Deployment}

\subsection{Supported Formats}

\begin{itemize}
    \item \textbf{SafeTensors}: BF16 full-precision weights (28.6 GB)
    \item \textbf{GGUF}: Q4\_K\_M (8.1 GB), Q5\_K\_M (9.8 GB), Q8\_0 (14.7 GB)
    \item \textbf{MLX}: 4-bit and 8-bit for Apple Silicon
    \item \textbf{vLLM}: Continuous batching with PagedAttention
    \item \textbf{TensorRT-LLM}: Optimized for NVIDIA Hopper/Ada
\end{itemize}

\subsection{Hardware Requirements}

\begin{table}[H]
\centering
\begin{tabular}{llll}
\toprule
\textbf{Format} & \textbf{VRAM} & \textbf{Context} & \textbf{Hardware} \\
\midrule
BF16 (full)     & 32 GB & 128K  & A100 40GB $\times 1$ \\
BF16 (full)     & 80 GB & 1M    & H100 80GB $\times 1$ \\
Q4\_K\_M        & 10 GB & 32K   & RTX 3080 10GB \\
Q4\_K\_M        & 16 GB & 128K  & RTX 4080 Super \\
MLX 4-bit       & 12 GB & 64K   & M3 Pro (18GB) \\
MLX 4-bit       & 24 GB & 512K  & M4 Max (64GB) \\
\bottomrule
\end{tabular}
\caption{Deployment Hardware Requirements}
\end{table}

\subsection{API Usage}

\begin{lstlisting}[language=Python, caption=Zen4 with Native Tool Call]
from hanzo import HanzoAI

client = HanzoAI()

tools = [{
    "type": "function",
    "function": {
        "name": "get_weather",
        "description": "Get current weather for a city",
        "parameters": {
            "type": "object",
            "properties": {
                "city": {"type": "string"},
                "units": {"type": "string", "enum": ["celsius", "fahrenheit"]}
            },
            "required": ["city"]
        }
    }
}]

response = client.chat.completions.create(
    model="zen4-14b-instruct",
    messages=[{"role": "user", "content": "What is the weather in Tokyo?"}],
    tools=tools,
    tool_choice="auto"
)
\end{lstlisting}

\section{Safety and Alignment}

\subsection{Evaluation}

Zen4 was evaluated on the standard Hanzo safety harness covering:

\begin{table}[H]
\centering
\begin{tabular}{lcc}
\toprule
\textbf{Safety Metric} & \textbf{Zen3} & \textbf{Zen4} \\
\midrule
Harmful content refusal rate   & 94.1\% & 97.3\% \\
Over-refusal rate (benign)     & 8.4\%  & 4.1\%  \\
Jailbreak resistance (AdvBench)& 91.2\% & 96.4\% \\
Factual accuracy (TruthfulQA)  & 71.3\% & 78.9\% \\
\bottomrule
\end{tabular}
\caption{Safety and Alignment Metrics}
\end{table}

\subsection{Limitations}

\begin{itemize}
    \item Long-context recall degrades for documents with high token density and many
    distinct facts; precision on 1M context tasks averages 91.7\% but drops to
    approximately 82\% on adversarial multi-needle tasks.
    \item Tool-call argument generation occasionally hallucinate valid-looking but incorrect
    values for arguments not strongly constrained by the schema.
    \item Mathematical reasoning on competition-level problems (AIME, USAMO) benefits
    from extended thinking; Zen4-Thinking is recommended for those tasks.
\end{itemize}

\section{Conclusion}

Zen4 establishes a new performance ceiling for the 14B dense parameter class through
architectural innovation rather than scale. The Hierarchical Hybrid Attention mechanism
resolves the context-length memory wall, native tool pretraining elevates function-call
reliability from a fine-tuning artifact to a first-class capability, and the fourth
generation of Zen MoDE distillation closes the gap to models more than twice Zen4's size.
The resulting 40\% efficiency gain over Zen3 makes Zen4 the recommended general-purpose
deployment for applications requiring frontier reasoning, long-document understanding,
or reliable agentic tool use within a single-GPU budget.

\section*{Acknowledgments}

We thank the research teams at Hanzo AI and Zoo Labs Foundation, our infrastructure
partners, and the open-source community whose tooling underpins evaluation and deployment.

\bibliographystyle{plain}
\bibliography{references}

\appendix

\section{Model Card}

\begin{table}[H]
\centering
\begin{tabular}{ll}
\toprule
\textbf{Field} & \textbf{Value} \\
\midrule
Model Name      & Zen4 \\
Version         & v2026.01 \\
Release Date    & January 2026 \\
Parameters      & 14.3B \\
Context Length  & 1,048,576 tokens \\
License         & Apache 2.0 \\
Repository      & \href{https://huggingface.co/zenlm/zen4-14b-instruct}{huggingface.co/zenlm/zen4-14b-instruct} \\
Documentation   & \href{https://papers.zenlm.org/zen4}{papers.zenlm.org/zen4} \\
Contact         & research@hanzo.ai \\
\bottomrule
\end{tabular}
\caption{Zen4 Model Card}
\end{table}

\section{Benchmark Methodology}

All evaluations use greedy decoding (temperature 0) unless otherwise specified. MMLU uses
5-shot chain-of-thought prompting. MATH uses 4-shot prompting with the standard MATH
evaluation harness. HumanEval uses 0-shot with pass@1 metric. GPQA uses 0-shot with
the Diamond subset. IFEval uses the prompt-level strict accuracy metric. Context-length
evaluations use positions sampled uniformly across the full context window.

\end{document}
